%%%%%%%%%%%%%%%%%%%%%%%%%%%%%%%%%%%%%%%%%%%%%%%%%%%%%%%%%%%%%
%%  PLANTILLA IEEE UMNG - INFORME DE PRACTICA DE LABORATORIO %%
%%  PROGRAMA DE INGENIERIA MECATRONICA - BOGOTA, COLOMBIA     %%
%%  Autor plantilla: Sergio Andres Chaparro Moreno            %%
%%  Version 1.0 - Julio 2018                                  %%
%%%%%%%%%%%%%%%%%%%%%%%%%%%%%%%%%%%%%%%%%%%%%%%%%%%%%%%%%%%%%
%%  CONTENIDO: Robot movil inteligente con ESP32-S3          %%
%%  NOTA: esqueleto. Las secciones marcadas [TODO] no tienen  %%
%%  contenido real todavia. NO inventar resultados.           %%
%%%%%%%%%%%%%%%%%%%%%%%%%%%%%%%%%%%%%%%%%%%%%%%%%%%%%%%%%%%%%
\documentclass[journal]{IEEEtran}
\IEEEoverridecommandlockouts

%%%%%%%%%%%%%%%%%%%%%%%%%%%%
%%%%%%%% PAQUETES %%%%%%%%%%
%%%%%%%%%%%%%%%%%%%%%%%%%%%%
\usepackage{fancyhdr}
\usepackage{graphicx}
\usepackage[spanish, es-tabla]{babel}
\usepackage[utf8]{inputenc}
\usepackage{color}
\usepackage{hyperref}
\usepackage{wrapfig}
\usepackage{array}
\usepackage{multirow}
\usepackage{adjustbox}
\usepackage{nccmath}
\usepackage{subfigure}
\usepackage{amsfonts,latexsym}
\usepackage{enumerate}
\usepackage{booktabs}
\usepackage{float}
\usepackage{threeparttable}
\usepackage{array,colortbl}
\usepackage{ifpdf}
\usepackage{rotating}
\usepackage{cite}
%% stfloats es opcional (flotantes a dos columnas en el pie). Si no esta
%% instalado en tu distribucion, el informe compila igualmente.
\IfFileExists{stfloats.sty}{\usepackage{stfloats}}{}
\usepackage{url}
\usepackage{listings}
\usepackage{amsmath}
\usepackage{makecell}
\setcellgapes{2pt}

\newcolumntype{P}[1]{>{\centering\arraybackslash}p{#1}}
\newcommand{\tabitem}{~~\llap{\textbullet}~~}
\newcommand{\ctt}{\centering\scriptsize\textbf}
\newcommand{\dtt}{\scriptsize\textbf}
\renewcommand\IEEEkeywordsname{Palabras clave}

\hyphenation{op-tical net-works semi-conduc-tor}
\graphicspath{ {figs/} }

%%%%%%%%%%%%%%%%%%%%%%%%%%%%%%%%%%%%%%%%%%%%%%%%%%%%%%%%%%%%%
%%%%%%%%%%%%% ENCABEZADO DE PAGINAS TIPO UMNG %%%%%%%%%%%%%%%
%%%%%%%%%%%%%%%%%%%%%%%%%%%%%%%%%%%%%%%%%%%%%%%%%%%%%%%%%%%%%
\newcommand{\MYhead}{\smash{\scriptsize
\hfil\parbox[t][\height][t]{\textwidth}{\centering
\begin{picture}(0,0) \put(-0,-17){\IfFileExists{LogoUMNG}{\includegraphics[width=33mm]{LogoUMNG}}{}} \end{picture} \hspace{6.4cm}
INFORME DE PRACTICA DE LABORATORIO \hspace{5.15cm} Version 1.0\\
\hspace{6.45cm} PROGRAMA DE INGENIERIA MECATRONICA \hspace{4.65cm} Periodo 2026-2\\
\underline{\hspace{ \textwidth}}}\hfil\hbox{}}}
\makeatletter
\def\ps@headings{%
\def\@oddhead{\MYhead}%
\def\@evenhead{\MYhead}}%
\def\ps@IEEEtitlepagestyle{%
\def\@oddhead{\MYhead}%
\def\@evenhead{\MYhead}}%
\makeatother
\pagestyle{headings}
\addtolength{\footskip}{0\baselineskip}
\addtolength{\textheight}{-1\baselineskip}

%%%%%%%%%%%%%%%%%%%%%%%%%%%%%%%%
%%%%% INICIO DEL DOCUMENTO %%%%%
%%%%%%%%%%%%%%%%%%%%%%%%%%%%%%%%
\begin{document}

%%%%%%%%%%%%%%%%%%%%%%%%%%%%
%%% TITULO DEL DOCUMENTO %%%
%%%%%%%%%%%%%%%%%%%%%%%%%%%%
\title{Robot movil inteligente con ESP32-S3: control manual y por voz,
       brujula digital, deteccion de colores y evasion de obstaculos}

%%%%%%%%%%%%%%%%%%%%%%%%%%%%
%%%%%%%%% AUTORES %%%%%%%%%
%%%%%%%%%%%%%%%%%%%%%%%%%%%%
%% [TODO] Completar con los autores reales y el profesor. NO inventar nombres.
\author{[NOMBRE DEL AUTOR 1],~y~[NOMBRE DEL AUTOR 2]\\
        \textit{\{correo1,~correo2\}@unimilitar.edu.co}\\
        Profesor:~[NOMBRE DEL PROFESOR]%
\thanks{El presente documento corresponde a un informe de practica de laboratorio presentado en la Universidad Militar Nueva Granada durante el periodo 2026-2.}}

\maketitle

%%%%%%%%%%%%%%%%%%%%%
%%%%%% RESUMEN %%%%%%
%%%%%%%%%%%%%%%%%%%%%
\begin{abstract}
[TODO] Resumen de 150--250 palabras: que se diseno, con que hardware, que
subsistemas integra y que se valido experimentalmente. No incluir resultados
que no se hayan medido. No citar referencias en el resumen.
\end{abstract}

\begin{IEEEkeywords}
robot movil, ESP32-S3, brujula digital, sensor de color, evasion de obstaculos,
control por voz, navegacion autonoma, IoT.
\end{IEEEkeywords}

%%%%%%%%%%%%%%%%%%%%%%%%%%%%%%%%%%%%%
%%% SECCIONES (esqueleto) %%%%%%%%%%
%%%%%%%%%%%%%%%%%%%%%%%%%%%%%%%%%%%%%

\section{Introduccion y planteamiento del problema}
[TODO] Contexto, motivacion y problema. Citar fuentes verificables.

\section{Objetivos}
\subsection{Objetivo general}
[TODO]
\subsection{Objetivos especificos}
\begin{enumerate}[1.]
\item [TODO]
\end{enumerate}

\section{Requisitos funcionales y no funcionales}
[TODO] Ver \texttt{docs/requisitos.md}. Resumir RF-01\dots RF-09 y los RNF.

\section{Marco teorico}
[TODO] Navegacion diferencial, magnetometria y rumbo, espacio de color
(RGB/XYZ), sensores ultrasonicos, comunicacion inalambrica, maquinas de estados.

\section{Arquitectura del sistema}
[TODO] Ver \texttt{docs/arquitectura.md}. Incluir diagrama de bloques.

\section{Diseno mecanico}
[TODO] Chasis de acrilico, tren de traccion (2 ruedas + caster), cotas y CAD.
Pendiente de definir la pista y las cotas finales.

\section{Diseno electronico y alimentacion}
[TODO] ESP32-S3, TB6612FNG, sensores I2C y ultrasonico, banco Li-ion 3S + BMS
+ buck a 6 V/5 V, GND comun. Ver \texttt{hardware/pinout.md}.

\section{Seleccion y caracterizacion de sensores}
[TODO] QMC5883L, TCS34725, HC-SR04. Calibraciones y limites reales.

\section{Desarrollo del firmware}
[TODO] Modulos por subsistema (ver \texttt{firmware/}), temporizacion no
bloqueante, maquina de estados.

\section{Protocolo de comunicacion y aplicacion web}
[TODO] Protocolo JSON versionado, WebSocket, app web de control y telemetria.

\section{Diseno del sistema autonomo}
[TODO] Control de rumbo y evasion de obstaculos. Definir el error angular como:

\begin{equation}
\mathrm{error} = ((rumbo_{deseado} - rumbo_{actual} + 540) \bmod 360) - 180
\label{eq:error_angular}
\end{equation}

\section{Implementacion del chatbot (JARVIS)}
[TODO] Asistente por reglas + telemetria real; opcional LLM externo (clave
fuera del repositorio).

\section{Metodologia de pruebas}
[TODO] Pruebas por modulo y por nivel (1: 4 obst./3 colores; 2: 6/4; 3: 8/6).
Ver \texttt{docs/trazabilidad.md}.

\section{Resultados experimentales}
[TODO] Solo resultados reales medidos. Tablas y graficas.

\section{Analisis de errores y limitaciones}
[TODO]

\section{Conclusiones y trabajo futuro}
[TODO]

%%%%%%%%%%%%%%%%%%%%%%%%%%
%%%%%% BIBLIOGRAFIA %%%%%%
%%%%%%%%%%%%%%%%%%%%%%%%%%
\begin{thebibliography}{9}
%% [TODO] Referencias verificables en formato IEEE. Ejemplos de partida:
\bibitem{espressif} Espressif Systems, \emph{ESP32-S3 Technical Reference Manual}.
[En linea]. Disponible en: \url{https://www.espressif.com/}
\bibitem{qst} QST, \emph{QMC5883L Datasheet}.
\bibitem{ams} ams, \emph{TCS34725 Color Sensor Datasheet}.
\bibitem{toshiba} Toshiba, \emph{TB6612FNG Dual Motor Driver Datasheet}.
\end{thebibliography}

%% [TODO] Enlace al repositorio de GitHub:
\section*{Repositorio}
\url{https://github.com/[USUARIO]/PROYECTO_LAB}

\end{document}
